\chapter{交流—交流变流电路}

把一种形式的交流电变成另一种形式的交流电（改变电压、电流，或改变频率、相数）称为\textbf{交流—交流变流}。只改变电压或电流而不改变频率的，称为\textbf{交流电力控制电路}，其中输出电压可连续调节的称为\textbf{交流调压电路}；既改变电压又改变频率的，称为\textbf{交—交变频电路}（周波变流器）。

% ============================================================
\section{交流调压电路}

交流调压电路用两只晶闸管反并联（或用一只双向晶闸管）串在负载回路中，在每个半周内控制其导通起始时刻，从而改变负载上电压的有效值。与相控整流不同，交流调压的负载直接接在交流电源上，\textbf{没有整流环节}。

\subsection{单相交流调压电路}

\subsubsection*{1. 电阻负载}

设电源电压 $u_1=\sqrt{2}U_1\sin\omega t$，晶闸管在 $\alpha$ 处触发导通，到 $\pi$ 处电流过零关断；负半周由反并联的管子重复。负载电压为缺角的正弦波。

\keyformula{U_o=U_1\sqrt{\frac{1}{\pi}\left(\pi-\alpha+\frac{\sin 2\alpha}{2}\right)}=U_1\sqrt{\frac{\pi-\alpha}{\pi}+\frac{\sin2\alpha}{2\pi}}}

$\alpha=0$ 时 $U_o=U_1$（全导通）；$\alpha=\pi$ 时 $U_o=0$。故\textbf{移相范围为 $0\sim\pi$}。负载电流 $I_o=U_o/R$，负载消耗的功率 $P=U_o^2/R$；由于负载为纯电阻，\textbf{功率因数就等于电压比}：

\keyformula{\lambda=\frac{P}{S}=\frac{U_oI_o}{U_1I_o}=\frac{U_o}{U_1}=\sqrt{\frac{\pi-\alpha}{\pi}+\frac{\sin2\alpha}{2\pi}}}

\subsubsection*{2. 阻感负载}

阻感负载时电流滞后于电压，晶闸管的导通角 $\theta$ 由 $\alpha$ 与负载阻抗角 $\varphi=\arctan(\omega L/R)$ 共同决定，满足超越方程

\keyformula{\sin(\alpha+\theta-\varphi)=\sin(\alpha-\varphi)\,\mathrm{e}^{-\frac{\theta}{\omega\tau}},\qquad \tau=\frac{L}{R}}

\corepoint{阻感负载的交流调压有一个\textbf{极重要的结论}：若 $\alpha<\varphi$，则触发脉冲到来时电流尚未过零（上一半周的电流还在流），晶闸管无法关断，电路实际上\textbf{失去控制作用}，输出电压等于电源电压（不可控）；因此\textbf{必须保证 $\alpha\ge\varphi$}，即最小触发角为 $\varphi$。若要既能调压又能用窄脉冲触发，须采用\textbf{宽脉冲}或\textbf{脉冲列}（因电流可能延续到脉冲消失之后）。}

\subsubsection*{3. 晶闸管电流定额}

单相交流调压中每只晶闸管流过半个周期的负载电流，其有效值 $I_T=I_o/\sqrt{2}$。选用器件时须按此有效值折算到通态平均电流（见 2.3.3 节的 1.57 系数），并注意\textbf{阻感负载时电流波形非正弦}，不能用简单的正弦公式。

\subsection{三相交流调压电路}

三相交流调压电路按其联结形式有多种，各有不同的适用场合与限制。

\begin{table}[htbp]
\centering
\caption{三相交流调压电路的联结形式}
\label{tab:6-1}
\begin{tabular}{p{3.6cm}p{4.6cm}p{6.2cm}}
\toprule
\textbf{联结形式} & \textbf{构成与特点} & \textbf{使用注意} \\
\midrule
星形联结带中性线 & 三个单相调压电路分别接在三相四线制上 & 三次谐波电流能经中性线流通，中性线电流大；只能用于小容量 \\
星形联结不带中性线 & 三只双向晶闸管（或三对反并联晶闸管）接成星形，负载星形接 & 要求 $\alpha\ge\varphi$；\textbf{须用宽脉冲或双窄脉冲}，否则某相可能因无触发而不能形成回路；三次谐波无通路 \\
三角形联结 & 三只双向晶闸管接在负载三角形内部 & 线电流中三次谐波不能流通，但相电流中有 \\
支路控制三角形联结 & 晶闸管控制三角形负载的支路电流（TCR，晶闸管控制电抗器） & 用于静止无功补偿，连续调节感性无功 \\
中点控制三角形联结 & 三相负载接成三角形，用三只晶闸管控制其中点 & 用得较少 \\
\bottomrule
\end{tabular}
\end{table}

\corepoint{三相交流调压最常考“星形不带中性线”：由于没有中性线，\textbf{三次及三的倍数次谐波电流没有通路}，故线电流中不含 $3k$ 次谐波（只含 $6k\pm1$ 次）。另外\textbf{窄脉冲不行}——因为每相电流要靠另两相构成回路，脉冲必须足够宽（或用双窄脉冲）才能保证成对导通。}

% ============================================================
\section{其他交流电力控制电路}

\subsection{交流调功电路}

交流调功电路与交流调压电路\textbf{电路形式完全相同}，区别只在控制方式：它不是在每个周期内控制导通相位，而是\textbf{以交流电的周期为单位}控制通断——接通几个周期、再断开几个周期，通过改变通断周期数的比值来调节输出功率。这种控制方式也称\textbf{周波数控制}或“通断控制”。

设控制周期为 $N$ 个电源周期，其中接通 $n$ 个周期，则输出功率

\keyformula{P=\frac{n}{N}P_{\max},\qquad P_{\max}=\frac{U_1^2}{R}}

\begin{itemize}
  \item 优点：晶闸管在电压过零时通断，\textbf{不产生相位控制的谐波}，对电网的干扰小；
  \item 缺点：输出功率是阶梯式调节，且通断周期较长时会引起负载温度脉动（对温度惯性大的负载如电炉、温控系统正合适）；
  \item 适用场合：电炉温度控制、空调压缩机、电热水器等\textbf{惯性大、对谐波敏感}的负载。
\end{itemize}

\subsection{交流电力电子开关}

把晶闸管（或双向晶闸管）反并联接在负载回路中，只作“通/断”用而不调压，就是\textbf{交流电力电子开关}（固态开关，SSR）。

\begin{itemize}
  \item 与机械式接触器相比：\textbf{无触点、无电弧、动作快、寿命长、无噪声}，可实现频繁通断；
  \item 与交流调压相比：只需判断“通”或“断”，控制简单，通态时晶闸管全导通、损耗小；
  \item 常与补偿电容配合构成\textbf{晶闸管投切电容器（TSC）}，用于无功补偿——在电容电压与电网电压相等的时刻投入，可避免冲击电流。
\end{itemize}

% ============================================================
\section{交—交变频电路}

\textbf{交—交变频电路}（周波变流器，Cycloconverter）不经过中间直流环节，直接把一种频率的交流电变成另一种频率的交流电，因此也称“直接变频”。

\subsection{单相交—交变频电路}

\subsubsection*{1. 电路构成和基本工作原理}

单相交交变频电路由\textbf{两组反并联的相控整流电路}（正组 P 与反组 N）构成：正组工作时输出正半周电流，反组工作时输出负半周电流。改变两组整流的控制角 $\alpha$，就可使输出电压平均值按正弦规律变化。

\begin{itemize}
  \item 输出电压的幅值由 $\alpha$ 决定（$\alpha$ 在 $0\sim\pi/2$ 之间变化时 $U_d$ 由正的最大值变到零）；
  \item 输出电压的频率由 $\alpha$ 变化的速率决定——$\alpha$ 按正弦规律从 $\pi/2$ 减小到某个最小值再回到 $\pi/2$，输出就完成一个正弦半周；
  \item 这种控制方法称为\textbf{余弦交点法}：把期望输出的正弦波与各段线电压的余弦曲线求交，交点时刻即为触发时刻，从而使输出平均电压逼近正弦。
\end{itemize}

\subsubsection*{2. 输出频率的上限}

由于交交变频靠“拼”电源电压的片段来合成正弦波，一个输出周期内必须包含足够多的电源电压片段才能保证波形质量，因此\textbf{输出频率不能高于输入频率}：

\keyformula{f_o\le\left(\frac{1}{3}\sim\frac{1}{2}\right)f_i}

即电网 50 Hz 时，交交变频的输出频率上限约为 20 Hz 左右——这使它特别适合\textbf{低速大功率}传动（如轧机、矿井提升机、水泥磨机）。

\subsubsection*{3. 无环流与有环流控制}

正、反两组反并联，如果两组同时导通，会在两组之间产生\textbf{环流}（不经过负载的短路电流）。处理方式有两种：

\begin{itemize}
  \item \textbf{无环流控制}：任何时刻只允许一组工作，另一组封锁脉冲；切换时须先封锁、等电流过零并留死区，再开放另一组，因此\textbf{切换有死区、输出波形有畸变}，但不需要环流电抗器，效率高——实际多用这种。
  \item \textbf{有环流控制}：两组同时加触发脉冲，用\textbf{环流电抗器}限制环流；切换平滑、无死区，但需笨重的电抗器、损耗大。
\end{itemize}

\subsection{三相交—交变频电路}

三相交交变频电路由三组（每相一组）单相交交变频电路构成，按输出绕组连接方式分为\textbf{公共交流母线进线方式}与\textbf{输出星形联结方式}。

\begin{itemize}
  \item \textbf{公共交流母线进线方式}：三组单相变频电路的输入端接在同一交流母线上，输出端分别接电动机的三相绕组（绕组须分开、各有 6 个端子），三组相互独立，控制简单。
  \item \textbf{输出星形联结方式}：三组输出接成星形，电动机三相绕组的一端连成中性点。由于中性点不引出，\textbf{三组的输出电流之和为零}，可实现三相电流的自动平衡；但每组输出端对地电压较高，绝缘要求高。
\end{itemize}

三相交交变频的特点：\textbf{输出频率低（$\le f_i/3$）、功率大、效率高、无需中间直流环节}，但\textbf{功率因数低}（因相控整流在深控时 $\cos\varphi$ 很小）、\textbf{输入电流谐波大}、输出波形含较多谐波。

% ============================================================
\section{矩阵式变频电路}

\textbf{矩阵式变频电路}（Matrix Converter）把三相交流输入直接变成三相交流输出，中间不经任何储能元件，只用 9 个\textbf{双向开关}（每个由两只反并联的全控器件构成）按一定规律通断，把输入的三相电压“切片”后拼成输出三相电压。

\begin{itemize}
  \item \textbf{理论上可输出任意频率}（不受 $f_i/3$ 的限制），且输入电流可控制为正弦、功率因数可调（可做到 $\lambda\approx1$）；
  \item 无中间储能元件（无大电容、大电感），\textbf{功率密度高、体积小、寿命长}；
  \item 约束条件：（1） 任一输出相不能与两个输入相同时接通（否则输入线电压短路）；（2） 任一输入相不能同时断开所有输出（否则感性负载电流无通路）；这就是“\textbf{双向开关的安全换流}要求”，实际须采用四步换流或两步换流逻辑；
  \item 缺点：器件数量多（9 个双向开关 = 18 个全控器件）、驱动与保护复杂、电压传输比有限（理论上最大约 0.866），故目前主要用于中小功率、要求高功率密度的场合。
\end{itemize}

\corepoint{本章复习抓“三条线”：\textbf{（1） 交流调压}——电阻负载 $U_o=U_1\sqrt{\frac{\pi-\alpha}{\pi}+\frac{\sin2\alpha}{2\pi}}$、$\lambda=U_o/U_1$，阻感负载必须 $\alpha\ge\varphi$ 否则失控；\textbf{（2） 交流调功}——电路同调压但按周期通断，$P=(n/N)P_{\max}$，谐波小；\textbf{（3） 交交变频}——两组反并联相控整流 + 余弦交点法，输出频率 $\le f_i/3$，适合低速大功率，功率因数低。矩阵式变频则是“无储能直接变换”的方向性电路。}
